\documentclass[10pt,a4paper]{amsart}
\usepackage[margin=19mm]{geometry}
\usepackage{amsmath,amssymb,mathtools}
\usepackage[scr=rsfs,cal=euler]{mathalfa}
\usepackage{xfrac}
\usepackage[unicode,hidelinks,bookmarks=false]{hyperref}
\newcommand{\ie}{i.e.\ }
\newcommand{\cf}{cf.\ }
\newcommand{\resp}{respectively\ }
\newcommand{\Subsection}{\subsection}
\providecommand{\nobreakdash}{\nobreak\hbox{-}}
\setlength{\emergencystretch}{2em}
\multlinegap0pt
\usepackage{tikz}
\usetikzlibrary{cd}
\usepackage{amssymb}
\usepackage{mathtools}
\usepackage[normalem]{ulem}
\usepackage{tikz}
\newtheorem{thm}{Theorem}
\newcommand{\BP}{{\mathbb{P}}}
\newcommand{\BQ}{{\mathbb{Q}}}
\newcommand{\CC}{{\mathcal C}}
\title{Full title of the article}
\author{Junliang Shen}
\date{Source: arXiv:2512.14995v2 (CC BY 4.0)}
\begin{document}
\maketitle

\noindent\emph{Source and licence.} Full title of the article. Version arXiv:2512.14995v2 (CC BY 4.0), distributed by arXiv under the Creative Commons Attribution 4.0 International licence, http://creativecommons.org/licenses/by/4.0/, as recorded in the arXiv licence metadata for this exact version and shown on the abstract page https://arxiv.org/abs/2512.14995v2. Full licence notice: \texttt{licenses/arxiv-2512.14995-CC-BY-4.0.txt}.\par\smallskip

\noindent\textit{Editorial scope.} Counted unit: the article's thm thm3.2 (source0.tex lines 952--954). The article's complete original proof of it is delivered with it (source0.tex lines 963--998, one \texttt{proof} environment of 36 lines, which the article places away from the statement). No mathematics is written, repaired or re-derived: the statement and the proof are the article's own lines, in the article's own notation, and no retained line is substituted or re-flowed.  Retained uncounted as the source-local context the proof needs: lines 393--398.  Nothing was omitted from any retained span, so the package carries no omission record.  The retained lines cite 1 works, whose entries are transcribed verbatim from the article's own bibliography source in the e-print and delivered with short editorial labels recorded in the item's reference mapping.  The retained lines contain the article's own diagram code, which the wrapper typesets with the standard tikz packages; no image file is delivered and the package declares no asset dependency.  Editorial formatting only, in addition: an AMS \texttt{amsart} preamble that restates the article's own declarations and macros used by the retained lines, namely the environments \texttt{thm} on separate counters, together with the standard packages those lines need.  The retained environments print in this document as Theorem 1 in order of appearance, whereas the article prints the same items with the numbering of its own full document.  The complete native source of the article as distributed in the arXiv e-print for this exact version is bundled unchanged as \texttt{sources/191139/source0.tex}.
\begin{thm}\label{thm1.1}
We have
\begin{equation}\label{DT_main}
R\pi_* \BQ_{\overline{J}_\CC} \simeq \bigoplus_{i=0}^{2g} \mathrm{IC}_B\left(  \wedge^i R^1p_* \BQ_{{\CC^\circ}}\right)[-i -\dim B] \in D^b_c(B).
    \end{equation}
\end{thm}

\begin{thm}\label{thm3.2}
    The decomposition of $R\pi_{k*} \BQ_{\CC^{[k]}}$ has full support for any $k\geq 1$.
\end{thm}

\begin{proof}[Proof sketch of Theorem \ref{thm3.2}]
The strategy of \cite{MY} is to use a descending induction on $k$. When $k \geq 2g-1$, the generalized Abel--Jacobi map $\mathrm{AJ}_b: \CC_b^{[k]} \to \overline{J}_{\CC_b}$ defined via the marked point $s(b) \in \CC_b$ is a projective bundle. Relatively over the base $B$, this gives a projective bundle
\[
\mathrm{AJ}_B: \CC^{[k]} \to \overline{J}_\CC
\]
defined via the section $s: B \to \CC $. Since the morphism $\pi_k:\CC^{[k]} \to B$ is the composition of $\mathrm{AJ}_B$ and $\pi: \overline{J}_\CC \to B$, we conclude for $k\geq 2g-1$ that $R\pi_{k*} \BQ_{\CC^{[k]}}$ has full support by Theorem \ref{thm1.1}.

To complete the proof, we need to show that any support of $\pi_{k}: \CC^{[k]} \to B$ also appears as a support of $\pi_{k+1}: \CC^{[k+1]} \to B$. Consider the incidence variety
\[
T_k \subset \left( \CC^{[k]} \times \BP^1 \right) \times_B \CC^{[k+1]}
\]
parameterizing $(Z, Z',a,b)$ where:
\begin{enumerate}
    \item[$\bullet$]  $b \in B$,
    \item[$\bullet$] $Z,Z' \subset \CC_b$ are of lengths $k, k+1$ respectively, \emph{i.e.} $[Z] \in \CC_b^{[k]}, [Z'] \in \CC_b^{[k+1]}$, such that $Z\subset Z'$,
    \item[$\bullet$] $a \in \BP^1$, 
    \item[$\bullet$] and the incidence condition is that the quotient of the ideal sheaves $I_{Z}/I_{Z'}$ is a skyscraper sheaf on $\CC_b$ over the point $a \in \BP^1$ via the finite map $\nu_b: \CC_b \to \BP^1$.
\end{enumerate} 

\medskip

Consider the natural projections
\begin{equation}\label{incidence}
    \begin{tikzcd}
T_k \arrow[r, ""] \arrow[d, ""]
& \CC^{[k+1]} \arrow[d, "\pi_{k+1}"]\\
\CC^{[k]} \times \BP^1 \arrow[r, "\pi'_k"] & B.
\end{tikzcd}
\end{equation}
Here $\pi'_k$ denotes the composition of the natural projection $\CC^{[k]} \times \BP^1 \to \CC^{[k]}$ and $\pi_k: \CC^{[k]} \to B$.
The incidence yields a relative correspondence over $B$,
\[
[T_k]: R\pi'_{k*} \BQ_{\CC^{[k]} \times \BP^1} \to R\pi_{k+1*} \BQ_{\CC^{[k+1]}}.
\]
A geometric study of the incidence (\ref{incidence}) shows that any simple summand of $R\pi'_{k*} \BQ_{\CC^{[k]} \times \BP^1}$ must appear as a simple summand of $ R\pi_{k+1*} \BQ_{\CC^{[k+1]}}$ via the correspondence $[T_k]$; see \cite[Page 10]{MY}. This completes the descending induction.
\end{proof}

\begin{thebibliography}{99}
\bibitem[MY]{MY} D. Maulik and Z. Yun, {\em Macdonald formula for curves with planar singularities,} J.Reine Angew. Math., 694 (2014), 27--48. %
\end{thebibliography}
\end{document}
